% ============================================================
% chapter3.tex — Chương 3: Phương pháp nghiên cứu
% ============================================================

\chapter{PHƯƠNG PHÁP NGHIÊN CỨU}
\label{chap:phuongphap}

Chương này mô tả chi tiết phương pháp nghiên cứu thông qua pipeline end-to-end, từ mô hình BIM đến dashboard trực quan hóa. Mỗi module trong pipeline được trình bày về đầu vào, quy trình xử lý và đầu ra.

% ---------------------------------------------------------------
\section{Pipeline tổng thể}
\label{sec:pipeline}
% ---------------------------------------------------------------

Pipeline nghiên cứu bao gồm các bước chính được minh họa trong Hình~\ref{fig:pipeline}.

\begin{figure}[htbp]
\centering
\resizebox{\textwidth}{!}{%
\begin{tikzpicture}[
  block/.style={
    rectangle,
    draw=NavyBlue,
    fill=NavyBlue!8,
    rounded corners=3pt,
    minimum width=2.6cm,
    minimum height=0.7cm,
    text centered,
    align=center,
    font=\scriptsize\bfseries,
    inner sep=3pt
  },
  fdsblock/.style={
    rectangle,
    draw=red!60!black,
    fill=red!8,
    rounded corners=3pt,
    minimum width=2.6cm,
    minimum height=0.7cm,
    text centered,
    align=center,
    font=\scriptsize\bfseries,
    inner sep=3pt
  },
  arrow/.style={
    -Stealth,
    semithick,
    NavyBlue,
  },
  redarrow/.style={
    -Stealth,
    semithick,
    red!60!black,
  }
]
  % === HÀNG 1: BIM & Graph ===
  \node[block] (revit) at (0, 5.0) {Revit / BIM};
  \node[block] (ifc) at (3.2, 5.0) {Export IFC};
  \node[block] (parse) at (6.4, 5.0) {IFC Parser\\(IfcOpenShell)};
  \node[block] (graph) at (9.6, 5.0) {Building Graph\\(NetworkX)};

  % === HÀNG 2: Simulation & Data ===
  \node[block] (scenario) at (0, 2.5) {Scenario Generator};
  
  \node[block] (cfast) at (3.2, 3.2) {CFAST Batch Sim};
  \node[fdsblock] (fds) at (3.2, 1.8) {FDS Reference Cases};
  
  \node[block] (post) at (6.4, 3.2) {Post-processing};
  \node[fdsblock] (compare) at (6.4, 1.8) {FDS Post-processing\\Aligned Sensors + Labels};
  \node[fdsblock] (fdsref) at (9.6, 1.1) {FDS Reference Set};
  
  \node[block] (iot) at (9.6, 2.5) {Virtual IoT\\Simulation};
  \node[block] (dataset) at (12.8, 2.5) {Graph-Time-Series\\Dataset};

  % === HÀNG 3: AI & Dashboard ===
  \node[block] (ai) at (0, -0.5) {AI Models:\\Baseline $\rightarrow$ GNN};
  \node[block] (eval) at (3.2, -0.5) {Evaluation:\\AI vs CFAST/FDS};
  \node[block] (api) at (6.4, -0.5) {Inference API};
  \node[block] (dash) at (9.6, -0.5) {Dashboard\\Prototype};

  % --- Đường nối Hàng 1 ---
  \draw[arrow] (revit) -- (ifc);
  \draw[arrow] (ifc) -- (parse);
  \draw[arrow] (parse) -- (graph);

  % --- Đường nối Hàng 1 -> Hàng 2 (Vòng ngoài qua rìa biên) ---
  \draw[arrow] (graph.east) -- +(0.4, 0) |- (11.3, 4.1) -- (-1.8, 4.1) |- (scenario.west);

  % --- Đường nối Hàng 2 ---
  \draw[arrow] (scenario.east) -- (cfast.west);
  \draw[redarrow] (scenario.east) -- (fds.west);
  
  \draw[arrow] (cfast) -- (post);
  \draw[redarrow] (fds) -- (compare);
  \draw[arrow] (post.east) -- (iot.west);
  \draw[arrow] (iot) -- (dataset);
  \draw[redarrow] (compare) -- (fdsref);

  % --- Đường nối Hàng 2 -> Hàng 3 (Vòng ngoài qua rìa biên) ---
  \draw[arrow] (dataset.east) -- +(0.5, 0) |- (14.6, 0.4) -- (-1.8, 0.4) |- (ai.west);

  % --- Đường nối Hàng 3 ---
  \draw[arrow] (ai) -- (eval);
  \draw[redarrow, dashed] (fdsref.south) -- (9.6, 0.1) -- (3.2, 0.1) -- (eval.north);
  \draw[arrow] (eval) -- (api);
  \draw[arrow] (api) -- (dash);

\end{tikzpicture}
}%
\caption{Pipeline tổng thể của đề tài}
\label{fig:pipeline}
\end{figure}

Bảng~\ref{tab:dataflow} trình bày data flow giữa các module.

\begin{table}[htbp]
\centering
\caption{Data flow giữa các module}
\label{tab:dataflow}
\footnotesize
\begin{tabularx}{\textwidth}{
  >{\raggedright\arraybackslash}p{0.21\textwidth}
  >{\raggedright\arraybackslash}p{0.29\textwidth}
  >{\raggedright\arraybackslash}X}
\toprule
\textbf{Module} & \textbf{Input} & \textbf{Output} \\
\midrule
IFC Parser        & File \texttt{.ifc}       & Bảng dữ liệu BIM thô \\
Graph Builder     & Bảng BIM                  & \texttt{building\_graph.json}, \texttt{node\_map.csv}, \texttt{edge\_map.csv} \\
Scenario Generator & Graph + factor configs   & Scenario configs \\
CFAST Runner      & Scenario configs           & Raw CFAST outputs \\
Post-processing   & Raw outputs                & Node time-series + labels \\
FDS Post-processing & Raw FDS outputs + sensor schema & Aligned FDS sensors + labels \\
CFAST--FDS Comparison & Aligned CFAST/FDS outputs & Reference comparison metrics \\
IoT Simulator     & CFAST node time-series     & Clean/noisy/faulty sensor streams \\
Dataset Builder   & Graph + scenario + sensors + labels & Model-ready CFAST dataset (PyG) \\
AI Trainer        & Dataset                    & Trained model \\
AI Evaluation     & AI predictions + CFAST test labels + FDS reference set & Metrics + ablation results \\
Dashboard API     & Graph + sensors + predictions & REST/WS endpoints \\
\bottomrule
\end{tabularx}
\end{table}

% ---------------------------------------------------------------
\section{Module BIM/IFC-to-Graph}
\label{sec:bimtograph}
% ---------------------------------------------------------------

\subsection{Quy trình xử lý}

Module BIM/IFC-to-Graph thực hiện chuyển đổi file IFC thành đồ thị tòa nhà theo quy trình:

\begin{enumerate}
  \item \textbf{Chuẩn bị trong Revit:} Kiểm tra và bổ sung Room/Space cho tất cả không gian (phòng, hành lang, cầu thang, sảnh). Export IFC với cài đặt bật room/space export.
  \item \textbf{Parse IFC:} Sử dụng \texttt{IfcOpenShell} để trích xuất \texttt{IfcSpace}, \texttt{IfcDoor}, \texttt{IfcWindow}, \texttt{IfcStair}, \texttt{IfcBuildingStorey} và \texttt{IfcRelSpaceBoundary}. Danh sách được xử lý theo từng loại entity, không ghép thành một chuỗi đầu vào duy nhất.
  \item \textbf{Xây dựng graph:} Sử dụng \texttt{NetworkX} để tạo đồ thị: mỗi \texttt{IfcSpace} thành node, mỗi cửa/lỗ mở nối hai space thành edge.
  \item \textbf{Gán thuộc tính:} Gán diện tích, thể tích, chiều cao, tầng, loại node cho mỗi node; gán loại edge, kích thước cửa cho mỗi edge.
  \item \textbf{Kiểm tra topology:} Kiểm tra graph có connected, có isolated nodes, có thiếu edge không.
\end{enumerate}

\subsection{Output}

\begin{itemize}
  \item \texttt{node\_map.csv}: Bảng ánh xạ \texttt{node\_id} $\leftrightarrow$ IFC \texttt{GlobalId}.
  \item \texttt{edge\_map.csv}: Bảng ánh xạ edge $\leftrightarrow$ IFC \texttt{GlobalId}.
  \item \texttt{building\_graph.json}: Topology + static features.
  \item \texttt{building\_graph.graphml}: File đồ thị cho visualization/debug.
  \item \texttt{graph\_check\_report.md}: Báo cáo lỗi topology và chỉnh sửa thủ công.
\end{itemize}

\subsection{Fallback}

Nếu IFC không đủ sạch, graph được xây dựng bán thủ công từ mặt bằng kiến trúc, room list và vị trí cửa. Mọi chỉnh sửa phải được ghi rõ trong tệp báo cáo kiểm tra topology.

% ---------------------------------------------------------------
\section{Thiết kế kịch bản cháy}
\label{sec:kichban}
% ---------------------------------------------------------------

\subsection{Nhóm biến và cấu hình thí nghiệm}

Để tránh trộn lẫn biến vật lý, cấu hình cảm biến, quy tắc tạo nhãn và quy tắc đánh giá, đề tài tổ chức cấu hình thành bốn nhóm độc lập.

\begin{table}[H]
\centering
\caption{Biến vật lý và cấu hình kịch bản P0}
\label{tab:factorPhysical}
\begin{tabularx}{\textwidth}{l l X}
\toprule
\textbf{Nhóm} & \textbf{Yếu tố} & \textbf{Biến / Tham số} \\
\midrule
Hình học  & Diện tích, thể tích, chiều cao phòng & \texttt{room\_area}, \texttt{room\_volume}, \texttt{ceiling\_height} \\
Hình học  & Kết nối corridor/stair/shaft          & \texttt{adjacency}, \texttt{is\_vertical} \\
Lỗ mở     & Diện tích, chiều cao cửa              & \texttt{opening\_area}, \texttt{opening\_height} \\
Lỗ mở     & Trạng thái cửa                        & \texttt{door\_state} \\
Cháy      & Vị trí cháy                           & \texttt{fire\_node}, \texttt{fire\_floor} \\
Cháy      & HRR, nhiên liệu, soot                 & \texttt{HRR(t)}, \texttt{growth\_class}, \texttt{soot\_yield} \\
Biên      & Nhiệt độ ban đầu / ngoài trời         & \texttt{T\_initial}, \texttt{T\_out} \\
Hệ thống  & Sprinkler/hút khói/tăng áp             & \texttt{sprinkler\_state}, \texttt{exhaust\_state} \\
\bottomrule
\end{tabularx}
\end{table}

\begin{table}[H]
\centering
\caption{Cấu hình IoT P0}
\label{tab:factorIoT}
\begin{tabularx}{\textwidth}{l X}
\toprule
\textbf{Nhóm} & \textbf{Biến / Tham số} \\
\midrule
Vị trí và loại cảm biến & \texttt{sensor\_node}, \texttt{sensor\_height}, \texttt{sensor\_type} \\
Lấy mẫu & \texttt{sensor\_sampling\_rate}, \texttt{sensor\_delay} \\
Suy giảm dữ liệu & \texttt{noise\_preset}, \texttt{missing\_mask}, \texttt{fault\_type}, \texttt{noise\_seed} \\
\bottomrule
\end{tabularx}
\end{table}

\begin{table}[H]
\centering
\caption{Cấu hình hậu xử lý và nhãn P0}
\label{tab:factorLabel}
\begin{tabularx}{\textwidth}{l X}
\toprule
\textbf{Nhóm} & \textbf{Biến / Tham số} \\
\midrule
Mô phỏng & \texttt{t\_end}, \texttt{dt\_out}, \texttt{scenario\_id} \\
Nhãn khói & \texttt{obscuration\_threshold}, \texttt{persistence\_steps}, \texttt{forecast\_horizon} \\
Thời gian đến & \texttt{arrival\_time\_abs}, \texttt{time\_to\_arrival}, \texttt{arrival\_valid\_mask} \\
\bottomrule
\end{tabularx}
\end{table}

\begin{table}[H]
\centering
\caption{Cấu hình split và thí nghiệm}
\label{tab:factorExperiment}
\begin{tabularx}{\textwidth}{l X}
\toprule
\textbf{Nhóm} & \textbf{Biến / Tham số} \\
\midrule
Phân chia & \texttt{base\_scenario\_id}, \texttt{split\_id}, \texttt{split\_mode} \\
Đánh giá & \texttt{model\_id}, \texttt{ablation\_id}, \texttt{evaluation\_seed} \\
Truy vết & \texttt{post\_processing\_version}, \texttt{label\_rule\_version} \\
\bottomrule
\end{tabularx}
\end{table}

\subsection{Lấy mẫu kịch bản}

Để tránh bùng nổ tổ hợp, nhóm sử dụng \textbf{scenario sampling có kiểm soát} thay vì full Cartesian product. Mỗi kịch bản vật lý được gán \texttt{base\_scenario\_id} duy nhất kèm metadata: \texttt{building\_id}, \texttt{fire\_node}, \texttt{fire\_floor}, \texttt{HRR\_config}, \texttt{door\_state\_config}, \texttt{system\_state\_config}, \texttt{boundary\_config}. Các biến thể nhiễu IoT dùng chung \texttt{base\_scenario\_id} và chỉ khác \texttt{noise\_seed}.

% ---------------------------------------------------------------
\section{Mô phỏng CFAST}
\label{sec:cfastmethod}
% ---------------------------------------------------------------

\subsection{Thiết lập mô phỏng}

Mỗi kịch bản cháy được chuyển đổi tự động thành file input CFAST dựa trên:
\begin{itemize}
  \item Thông tin khoang từ building graph (kích thước, cao độ).
  \item Thông tin lỗ thông gió từ edge (cửa, cửa sổ, trạng thái đóng/mở).
  \item Đường cong HRR và thông số nhiên liệu từ scenario config.
  \item Điều kiện biên từ boundary config.
\end{itemize}

\subsection{Batch simulation}

CFAST được chạy tự động cho tất cả kịch bản trong bộ scenario catalog. Script batch runner quản lý việc tạo input, chạy mô phỏng, thu thập output và log.

\subsection{Post-processing}

Output CFAST thô được xử lý thành:
\begin{itemize}
  \item \textbf{Node time-series:} Chuỗi thời gian cho mỗi node gồm nhiệt độ, chiều cao lớp khói, optical density và các đại lượng vật lý cần thiết.
  \item \textbf{Virtual sensor clean:} Giá trị tại vị trí/cao độ cảm biến theo một sensor schema cố định.
  \item \textbf{Labels P0:} Nhãn khói; thời điểm khói đến tuyệt đối; thời gian còn lại đến khi có khói; và mask xác định node có thời gian đến hợp lệ.
\end{itemize}

\subsection{Quy tắc tạo nhãn P0}

Đề tài sử dụng độ che khuất tại cảm biến khói $O_n(t)$, đơn vị \%/m, làm tiêu chí thống nhất cho nhãn P0. Theo CFAST User's Guide, giá trị setpoint mặc định của smoke detector là $23.93$~\%/m \cite{cfast_guide}. Đây là ngưỡng vận hành dữ liệu của nghiên cứu, không phải ngưỡng nguy hiểm pháp lý.

Với $\tau_O=23.93$~\%/m và số bước duy trì $m=2$, nhãn tại node $n$ được định nghĩa:
\begin{equation}
\texttt{smoke\_label}(n,t)=
\begin{cases}
1, & O_n(t-j\Delta t_{\mathrm{out}})\geq \tau_O,\quad j=0,\ldots,m-1,\\
0, & \text{ngược lại.}
\end{cases}
\end{equation}

Thời điểm khói đến tuyệt đối là thời điểm đầu tiên điều kiện trên được thỏa mãn:
\begin{equation}
\texttt{arrival\_time\_abs}(n)=
\min\{t:\texttt{smoke\_label}(n,t)=1\}.
\end{equation}
Nếu node không đạt ngưỡng trong thời gian mô phỏng, \texttt{arrival\_valid\_mask}$=0$ và không gán một thời gian giả. Tại thời điểm dự báo $t$, mục tiêu hồi quy là:
\begin{equation}
\texttt{time\_to\_arrival}(n,t)=
\max\left(\texttt{arrival\_time\_abs}(n)-t,0\right),
\end{equation}
và loss hồi quy chỉ được tính tại các phần tử có mask hợp lệ. Mục tiêu phân loại của mô hình là \texttt{smoke\_label}$(n,t+\Delta)$.

% ---------------------------------------------------------------
\section{Đối sánh CFAST--FDS}
\label{sec:fdsvalidation}
% ---------------------------------------------------------------

FDS được sử dụng để đối sánh chọn lọc trên một số kịch bản đại diện. Đây là so sánh giữa hai mô hình số, không phải kiểm chứng bằng thí nghiệm cháy thật.

\begin{table}[H]
\centering
\caption{Các case đối sánh FDS}
\label{tab:fdscases}
\begin{tabular}{l c l}
\toprule
\textbf{Case} & \textbf{Mức} & \textbf{Mục đích} \\
\midrule
FDS-01 & P0 & Phòng → hành lang \\
FDS-02 & P0 & Hành lang dài \\
FDS-03 & P0 & Phòng → hành lang → cầu thang \\
FDS-04 & P1 & Cửa mở so với cửa đóng \\
FDS-05 & P1 & Cửa sổ và gió \\
FDS-06 & P1 & Hút khói tắt/bật \\
FDS-07 & P2 & Tăng áp cầu thang tắt/bật \\
FDS-08 & P2 & Có/không sprinkler \\
\bottomrule
\end{tabular}
\end{table}

Đối với mỗi case, FDS sử dụng cùng hình học rút gọn, nguồn cháy, điều kiện biên, vị trí và cao độ cảm biến, bước kết xuất và quy tắc tạo nhãn như case CFAST tương ứng. Từ FDS, hệ thống tạo \texttt{virtual\_sensor\_clean}, các preset IoT và nhãn P0 bằng chính pipeline hậu xử lý dùng cho CFAST. Tập này chỉ dùng để đối sánh, không hòa vào tập huấn luyện CFAST.

Kết quả được sử dụng cho hai phép so sánh:
\begin{itemize}
  \item CFAST so với FDS trên các đại lượng đã căn chỉnh.
  \item AI so với nhãn FDS trong tập tham chiếu độc lập.
\end{itemize}

\subsection{Cơ chế quy đổi và metrics đối sánh}

FDS biểu diễn trường 3D, trong khi CFAST sử dụng hai vùng cho mỗi khoang. Vì vậy, báo cáo không so sánh trung bình toàn thể tích FDS trực tiếp với upper layer CFAST. Độ che khuất và nhiệt độ cảm biến được lấy tại cùng vị trí/cao độ. Khi cần so sánh đại lượng upper layer, dữ liệu FDS được trung bình trên các ô lưới nằm phía trên độ cao mặt phân cách lớp do CFAST xác định tại cùng thời điểm. Mọi chuỗi được nội suy về cùng \texttt{dt\_out} trước khi tính metric.

Sau khi đồng bộ hóa, các metric trong Bảng~\ref{tab:cfast_fds_metrics} được sử dụng để mô tả sai lệch. Không đặt ngưỡng ``đạt chuẩn pháp lý'' từ các metric này.

\begin{table}[H]
\centering
\caption{Các chỉ số đối sánh CFAST--FDS ở cấp node}
\label{tab:cfast_fds_metrics}
\begin{tabularx}{\textwidth}{l c X}
\toprule
\textbf{Đại lượng} & \textbf{Mức} & \textbf{Metrics} \\
\midrule
\texttt{smoke\_label} & P0 & F1, Recall, Precision \\
\texttt{arrival\_time\_abs} & P0 & MAE, RMSE trên các node có mask hợp lệ ở cả hai mô hình \\
Nhiệt độ và độ che khuất tại cảm biến & P0 & MAE, RMSE và tương quan chuỗi thời gian \\
\texttt{smoke\_path} & P1 & Path overlap trên thứ tự node bị ảnh hưởng \\
\bottomrule
\end{tabularx}
\end{table}

% ---------------------------------------------------------------
\section{Mô phỏng IoT giả lập}
\label{sec:iotmethod}
% ---------------------------------------------------------------

\subsection{Quy trình tạo dữ liệu IoT}

Từ output CFAST, hệ thống IoT giả lập tạo dữ liệu theo ba lớp:

\begin{enumerate}
  \item Lưu output gốc có truy vết theo node và thời gian → \texttt{physics\_truth}.
  \item Lấy mẫu tại vị trí/cao độ cảm biến cố định → \texttt{virtual\_sensor\_clean}.
  \item Thêm nhiễu Gauss, trễ, mất dữ liệu và lỗi cảm biến → \texttt{iot\_noisy\_stream}.
\end{enumerate}

Các case FDS P0 dùng cùng sensor schema và cùng quy tắc suy giảm để tạo tập tham chiếu tương thích, nhưng không được dùng để huấn luyện mô hình chính.

\subsection{Quy tắc chống label leakage}

Để đảm bảo tính công bằng của đánh giá AI:
\begin{itemize}
  \item Tại thời điểm $t$, AI chỉ được sử dụng building graph, cấu hình kịch bản đã biết và dữ liệu IoT từ $[t-k,t]$.
  \item Target phân loại là \texttt{smoke\_label}$(n,t+\Delta)$; cấm sử dụng output CFAST/FDS tại thời điểm sau $t$ làm input.
  \item Cấm sử dụng mọi label làm input, bao gồm nhãn khói, thời gian khói đến, thời gian còn lại và đường lan truyền.
  \item Toàn bộ biến thể \texttt{noise\_seed} của cùng \texttt{base\_scenario\_id} phải thuộc cùng một split.
\end{itemize}

% ---------------------------------------------------------------
\section{Xây dựng dataset}
\label{sec:dataset}
% ---------------------------------------------------------------

Dataset cuối cùng cho huấn luyện AI được tổ chức dưới dạng graph-time-series, bao gồm:
\begin{itemize}
  \item \texttt{building\_graph.json}: Topology đồ thị.
  \item \texttt{node\_static.parquet}: Đặc trưng tĩnh của node.
  \item \texttt{edge\_static.parquet}: Đặc trưng tĩnh của edge.
  \item \texttt{scenario\_static.parquet}: Nguồn cháy, HRR config và metadata kịch bản đã biết.
  \item \texttt{node\_dynamic.parquet}: Trạng thái động cấp node và hệ thống.
  \item \texttt{edge\_dynamic.parquet}: Trạng thái cửa/lỗ mở theo thời gian.
  \item \texttt{boundary\_dynamic.parquet}: Điều kiện biên theo thời gian.
  \item \texttt{iot\_noisy\_stream.parquet}: Dữ liệu cảm biến IoT (input chính).
  \item \texttt{labels.parquet}: Nhãn P0, thời gian đến tuyệt đối và mask.
  \item \texttt{splits/}: Phân chia train/val/test theo \texttt{base\_scenario\_id}.
  \item \texttt{pyg/*.pt}: Data objects cho PyTorch Geometric \cite{fey2019pyg}.
\end{itemize}

% ---------------------------------------------------------------
\section{Mô hình AI}
\label{sec:aimodel}
% ---------------------------------------------------------------

\subsection{Lộ trình phát triển mô hình}

Các họ mô hình được tổ chức từ đơn giản đến phức tạp. Phạm vi P0 bắt buộc triển khai tối thiểu ba baseline: rule-based, một baseline chuỗi thời gian không graph và một baseline graph; các kiến trúc còn lại là ứng viên lựa chọn trong từng họ.

\begin{table}[H]
\centering
\caption{Lộ trình phát triển mô hình AI}
\label{tab:modelroadmap}
\begin{tabularx}{\textwidth}{c l X}
\toprule
\textbf{Họ} & \textbf{Mô hình ứng viên} & \textbf{Mục đích} \\
\midrule
0 & Rule-based graph spread          & Baseline quy tắc bắt buộc \\
1 & Logistic Regression / RF / XGBoost / MLP & Baseline dạng bảng, tùy chọn \\
2 & LSTM / GRU / TCN                 & Baseline không graph bắt buộc \\
3 & GCN / GraphSAGE / GAT           & Baseline graph bắt buộc \\
4 & Temporal GNN                     & Kết hợp graph và lịch sử IoT \\
5 & Boundary-Aware Temporal GNN (BAT-GNN) & Mô hình chính \\
\bottomrule
\end{tabularx}
\end{table}

\subsection{Kiến trúc mô hình chính (BAT-GNN)}

Mô hình chính --- Boundary-Aware Temporal GNN --- thực hiện dự báo có điều kiện. Tại thời điểm $t$, mô hình chỉ nhận thông tin đã biết hoặc quan sát được đến $t$:

\begin{itemize}
  \item \textbf{Input:} Đồ thị $G(V,E)$; node/edge features tĩnh; cấu hình nguồn cháy và HRR đã biết; điều kiện biên, trạng thái cửa/hệ thống và lịch sử IoT trong $[t-k,t]$.
  \item \textbf{Encoders:}
  \begin{itemize}
    \item Graph Encoder: Encode cấu trúc không gian đồ thị.
    \item Temporal Encoder: Encode chuỗi thời gian IoT.
    \item Boundary Encoder: Encode nhiệt độ biên; gió chỉ bổ sung nếu triển khai factor P1.
    \item Fire Encoder: Encode thông tin nguồn cháy.
    \item System Encoder: Encode trạng thái hệ thống PCCC.
  \end{itemize}
  \item \textbf{Fusion:} Kết hợp các biểu diễn bằng concatenation hoặc attention.
  \item \textbf{Decoders:} Các đầu ra (multi-task) được phân tầng mức độ ưu tiên triển khai:
  \begin{itemize}
    \item \textbf{Mức P0 (Bắt buộc):} \texttt{smoke\_label}$(n,t+\Delta)$ và \texttt{time\_to\_arrival}$(n,t)$.
    \item \textbf{Mức P1 (Khuyến nghị/Tiêu chuẩn):} \texttt{hazard\_label} (nhãn nguy hiểm dựa trên nhiệt độ và nồng độ khí độc) và \texttt{unsafe\_time} (thời điểm bắt đầu không an toàn).
    \item \textbf{Mức P1/P2 (Mở rộng):} \texttt{next\_affected\_node} (node tiếp theo bị ảnh hưởng) và \texttt{smoke\_path} (đường lan truyền khói dự kiến).
  \end{itemize}
\end{itemize}

\subsection{Loss functions}

\begin{table}[H]
\centering
\caption{Loss function và mức độ ưu tiên cho từng task}
\label{tab:loss}
\begin{tabular}{l c l}
\toprule
\textbf{Output} & \textbf{Mức} & \textbf{Loss Function} \\
\midrule
\texttt{smoke\_label}        & P0 & BCE hoặc Focal Loss \\
\texttt{time\_to\_arrival}   & P0 & Masked MAE / Huber \\
\texttt{hazard\_label}       & P1 & BCE hoặc Focal Loss \\
\texttt{unsafe\_time}        & P1 & Masked MAE / Huber \\
\texttt{next\_affected\_node} & P1/P2 & Cross Entropy \\
Multi-task tổng hợp          & -- & Weighted sum \\
\bottomrule
\end{tabular}
\end{table}

% ---------------------------------------------------------------
\section{Đánh giá mô hình}
\label{sec:evaluation}
% ---------------------------------------------------------------

\subsection{Metrics}

\begin{table}[H]
\centering
\caption{Metrics đánh giá và mức độ ưu tiên theo task}
\label{tab:metrics}
\begin{tabularx}{\textwidth}{l c X}
\toprule
\textbf{Task} & \textbf{Mức} & \textbf{Metrics} \\
\midrule
Phân loại khói tại $t+\Delta$  & P0 & F1, Recall, Precision, AUC \\
Thời gian còn lại đến khi có khói & P0 & MAE, RMSE trên \texttt{arrival\_valid\_mask} \\
Phân loại nguy hiểm            & P1 & Recall, F1, AUC \\
Đường lan truyền khói          & P1/P2 & Path accuracy, next-node accuracy \\
Độ tin cậy dự báo (Calibration)& P1 & Brier score \\
Tốc độ dự đoán                 & P0 & Inference time \\
Độ bền vững (Robustness)       & P0 & Metric drop dưới điều kiện IoT bị nhiễu/lỗi \\
\bottomrule
\end{tabularx}
\end{table}

\begin{notebox}
Trong bài toán an toàn cháy, \textbf{Recall quan trọng hơn Accuracy} vì việc bỏ sót vùng nguy hiểm (false negative) nghiêm trọng hơn cảnh báo thừa (false positive).
\end{notebox}

\subsection{Tiêu chí thành công P0}

BAT-GNN được xem là đạt mục tiêu nghiên cứu P0 khi đồng thời thỏa mãn:
\begin{enumerate}
  \item F1 và Recall của \texttt{smoke\_label} trên test scenario cao hơn baseline không graph mạnh nhất.
  \item MAE \texttt{time\_to\_arrival} không lớn hơn baseline không graph mạnh nhất; nếu F1/Recall tăng nhưng MAE tăng thì phải báo cáo trade-off, không kết luận vượt trội toàn diện.
  \item Thời gian inference cho một scenario ngắn hơn thời gian chạy CFAST của chính scenario đó trên cùng cấu hình phần cứng báo cáo.
\end{enumerate}

\subsection{Phân chia dữ liệu}

\begin{table}[H]
\centering
\caption{Các chế độ phân chia dữ liệu}
\label{tab:splits}
\begin{tabularx}{\textwidth}{l X}
\toprule
\textbf{Split} & \textbf{Mục đích} \\
\midrule
Scenario split        & Đánh giá cơ bản \\
Hold-out fire node    & Kiểm tra tổng quát hóa với vị trí cháy mới \\
Hold-out fire floor   & Kiểm tra tổng quát hóa với tầng cháy mới \\
Noise robustness      & Đánh giá dưới điều kiện IoT suy giảm \\
FDS reference set     & Đối sánh chọn lọc với mô phỏng CFD chi tiết \\
Debug layout transfer & Kiểm tra pipeline chạy được trên layout thấp tầng; không dùng để tuyên bố tổng quát hóa \\
\bottomrule
\end{tabularx}
\end{table}

\subsection{Ablation study}

\begin{table}[H]
\centering
\caption{Thiết kế ablation study}
\label{tab:ablation}
\begin{tabularx}{\textwidth}{l X}
\toprule
\textbf{Thí nghiệm} & \textbf{Ý nghĩa} \\
\midrule
Không graph               & Baseline chuỗi thời gian với cùng lịch sử IoT và scenario features \\
Graph, không IoT          & Đóng góp của topology và cấu hình kịch bản \\
Graph + IoT               & Đóng góp của lịch sử cảm biến \\
Full without Boundary     & Đóng góp của điều kiện biên \\
Full without System       & Đóng góp của trạng thái hệ thống PCCC \\
Full model                & Tất cả feature hợp lệ tại thời điểm $t$ \\
\bottomrule
\end{tabularx}
\end{table}

% ---------------------------------------------------------------
\section{Dashboard prototype}
\label{sec:dashboard}
% ---------------------------------------------------------------

\subsection{Chức năng}

Dashboard prototype cung cấp:
\begin{itemize}
  \item Hiển thị đồ thị tòa nhà với filter theo tầng.
  \item Hiển thị trạng thái cảm biến ảo (khói, nhiệt, cửa, báo cháy).
  \item Hiển thị xác suất khói tại $t+\Delta$ và thời gian còn lại đến khi có khói tại mỗi node.
  \item Replay kịch bản cháy theo timeline.
  \item API inference AI cho cấu hình kịch bản và cửa sổ lịch sử IoT hợp lệ.
  \item So sánh AI với CFAST và, trên các case tham chiếu, với FDS.
\end{itemize}

\subsection{Công nghệ}

\begin{table}[H]
\centering
\caption{Tech stack dashboard}
\label{tab:techstack}
\begin{tabular}{l l}
\toprule
\textbf{Layer} & \textbf{Công nghệ} \\
\midrule
Backend          & FastAPI \\
AI inference     & PyTorch + FastAPI \\
Frontend         & React + TypeScript \\
Graph visualization & Cytoscape.js hoặc React Flow \\
Charts           & ECharts hoặc Recharts \\
Realtime         & WebSocket / SSE \\
\bottomrule
\end{tabular}
\end{table}

Dashboard \textbf{không} điều khiển hệ thống PCCC thật, không claim real-time CFD, và không hiển thị output không trace được về pipeline.
